\documentclass[10pt]{article}
\usepackage[margin=0.85in]{geometry}
\usepackage{amsmath,amssymb,amsthm,booktabs,hyperref}
\newtheorem{theorem}{Theorem}
\title{A counterexample to the Meijer-function zero-count conjecture}
\author{Principia Math in Progress}
\date{13 September 2026}
\begin{document}
\maketitle

\paragraph{Task 1. Resolve the zero-count conjecture.}
Conjecture 1 of Karp and Prilepkina \cite{KP} asserts that, for
\(a_i\ge0\) and \(\sum_i(b_i-a_i)>0\), the number of zeros of
\(G^{p,0}_{p,p}(x\mid b/a)\) in \((0,1)\) does not exceed the
number of zeros of \(v(t)=\sum_i(t^{a_i}-t^{b_i})\) there.
We refute it with strictly positive rational parameters.
The zero count of an explicit polynomial is certified by finite
integer arithmetic and Descartes' rule of signs.

\begin{theorem}
Let
\[
 a=\left(1,\frac{1999}{1000},\frac{3099}{1000},\frac{3447}{1000}\right),
 \qquad
 b=\left(5,\frac{999}{1000},\frac{2099}{1000},\frac{2447}{1000}\right).
\]
All parameters are positive and \(\sum_i(b_i-a_i)=1\).
The corresponding Meijer function has exactly three distinct
zeros in \((0,1)\), whereas \(v\) has exactly one, which is simple.
\end{theorem}

\begin{proof}
Put \(r=(-1,1099,1447)/1000\). The gamma recurrence gives
\[
 \prod_{i=1}^4\frac{\Gamma(s+a_i)}{\Gamma(s+b_i)}
 =\frac{\prod_{j=1}^3(s+1+r_j)}{(s+1)(s+2)(s+3)(s+4)}.
\]
By partial fractions, the Mellin transform on \((0,1)\) of
\(xF(x)/(6\cdot10^9)\) is this same rational function, where
\[
 F(x)=8859711253x^3-2991012759x^2+132891759x-1590253.
\]
Mellin inversion, or direct evaluation of the defining rational
Mellin--Barnes integrand, therefore gives
\[
 G^{4,0}_{4,4}(x\mid b/a)=\frac{xF(x)}{6\cdot10^9},
 \qquad 0<x<1.
\]
This identification involves no numerical evaluation of a special
function. The following four exact values have alternating signs:
\[
 F(0)=-1590253,\quad
 F(1/40)=\frac{69823293}{64000},\quad
 F(3/100)=-\frac{56199509269}{10^6},\quad
 F(3/10)=\frac{8298330221}{1000}.
\]
Thus the cubic has three distinct interior zeros, one in each
of the consecutive intervals determined by these points.
Its degree shows that there are no further zeros.

For the comparison function, put \(u=t^{1/1000}\). Directly
from the parameter vectors,
\begin{equation}\label{comparison}
 v(t)=u^{999}(1-u^{1000})q(u),\qquad
 q(u)=-1+u+u^{1001}-u^{1100}-u^{1448}+u^{2001}+u^{3001}.
\end{equation}
The prefactor is strictly positive for \(0<u<1\), so it remains
to count the zeros of \(q\) on that interval.

Here is the exact integer certificate. Let \(N=3001\), and for
integers \(0\le A<C\le B\) define
\[
 H_{A,C,B}(y)=(B(1+y))^N
    q\left(\frac{C+Ay}{B(1+y)}\right)\in\mathbb Z[y].
\]
The fractional-linear substitution maps \(y>0\) bijectively onto
\(A/B<u<C/B\). Multiplication by its positive prefactor preserves
zeros and their multiplicities. For each row of the table below,
the coefficient signs are listed in increasing degree; none is zero.
\[
\begin{array}{rrrcl}
\toprule
A&C&B&\text{interval}&\text{coefficient sign string of }H\\
\midrule
0&99&100&(0,99/100)&-^{3002}\\
990&995&1000&(99/100,199/200)&-^{3002}\\
995&997&1000&(199/200,997/1000)&-^{3002}\\
997&998&1000&(997/1000,499/500)&-^{3002}\\
998&999&1000&(499/500,999/1000)&-^{3002}\\
999&1000&1000&(999/1000,1)&+^{2587}-^{415}\\
\bottomrule
\end{array}
\]
Here \(-^m\) and \(+^m\) mean \(m\) consecutive minus and plus
signs, respectively. The table is reproducible with the integer
recurrence below. Two independently implemented coefficient
constructions agree on all \(6\cdot3002=18012\) coefficients.

Descartes' rule gives no roots in the first five intervals and
at most one root, counted with multiplicity, in the last.
Exact integer evaluation gives \(q<0\) at all six left endpoints
in the table and \(q(1)=1\). There are no roots at a subdivision
point. The intermediate value theorem supplies a root in the
last interval, so there is exactly one and it is simple.
Equation \eqref{comparison} proves the claimed count for \(v\),
and \(3>1\) disproves the conjecture.
\end{proof}

\paragraph{Reproducing the finite certificate.}
Write \(q(u)=\sum_i c_i u^i\), where its seven nonzero coefficients
are displayed in \eqref{comparison}. A summand of \(H\) is
\(c_i B^{N-i}(C+Ay)^i(1+y)^{N-i}\).
If its coefficients are \(h_k\), set
\[
 h_{-1}=0,\qquad h_0=c_i B^{N-i}C^i,\qquad
 \ell=iA+(N-i)C.
\]
Differentiating this product gives the recurrence
\begin{equation}\label{rec}
 h_{k+1}=\frac{(\ell-(A+C)k)h_k+A(N-k+1)h_{k-1}}{C(k+1)},
 \qquad 0\le k<N.
\end{equation}
Sum these coefficient lists over the seven values of \(i\).
Every division in \eqref{rec} is exact. The supplied
\texttt{verify.py} checks this, verifies the complete sign strings
and endpoint signs, and independently reconstructs \(H\) by
polynomial composition and reversal. Thus the root count covers
the entire real interval, including multiplicities; sampled plots
and floating-point root estimates are not used in the certificate.

\paragraph{Scope.}
The example lies strictly within the nonnegative-parameter and
positive-excess hypotheses. Taking all scale parameters equal
to one also refutes the more general Fox-function assertion
in Conjecture 3. This immediate consequence receives no additional
campaign credit.

\paragraph{Task 2. Find smaller counterexamples.}
Determine the least parameter count \(p\) for which the conjecture
fails, and seek examples with smaller parameter denominators.

\paragraph{Task 3. Find a replacement zero bound.}
Determine whether an alternative bound in terms of \(p\), the
parameter ordering, and the zeros of \(v\) holds for positive parameters.

\paragraph{Task 4. Describe counterexample regions.}
Identify explicit open regions of parameter space where the
Meijer function has three simple interior zeros and \(v\) has one.

\paragraph{Task 5. Study nonunit Fox scales.}
Construct counterexamples with unequal positive scale parameters
and analyze how the two zero counts change as those scales vary.

\begin{thebibliography}{1}
\bibitem{KP} D.~B.~Karp and E.~G.~Prilepkina,
\emph{Completely monotonic gamma ratio and infinitely divisible
H-function of Fox}, Comput. Methods Funct. Theory \textbf{16}
(2016), 135--153,
\href{https://doi.org/10.1007/s40315-015-0128-9}
{doi:10.1007/s40315-015-0128-9};
\href{https://arxiv.org/abs/1501.05388}{arXiv:1501.05388v2}.
Conjecture 1 is on p.~5 of the 18-page arXiv PDF and p.~140
(PDF page 8) of the inspected author-hosted journal offprint.
\end{thebibliography}
\end{document}
